\noindent Grupo:\rule{2cm}{0.4pt} 13 \rule{2cm}{0.4pt}

\vspace{0.6cm}

\noindent ¿Hicieron uso de herramientas de IA en su trabajo para este informe y/o presentación? \\
SI \rule{0.5cm}{0.4pt} X \rule{0.5cm}{0.4pt} \hspace{1cm} NO \rule{1.5cm}{0.4pt}

\vspace{0.7cm}

\noindent ¿Cuáles de las siguientes alternativas mejor describe el uso que hicieron de IA (puede ser más de una)

\vspace{0.4cm}

\noindent
\begin{tabular}{@{}p{0.02\textwidth} p{0.95\textwidth}@{}}
X & Apoyo a la redacción, ortografía, estructura gramatical, etc., de los documentos. \\
X & Apoyo al análisis de datos e información asociada al problema. \\
-  & Búsqueda y/o filtro de referencias, contenidos web, literatura, etc., relevante al proyecto. \\
X & Apoyo para el desarrollo de modelos asociados al proyecto. \\
X & Apoyo en el desarrollo de código computacional necesario para el proyecto \\
X & Apoyo en el análisis de resultados (por ejemplo, de distintas corridas computacionales, tendencias de predicciones, etc.) \\
x  & Apoyo para construir figuras, diagramas, etc. \\
- & Apoyo para obtener conclusiones a partir de los análisis de resultados. \\
X & Otro, especificar: Generación de un entorno Docker para visualización interactiva de datos financieros. \\
\end{tabular}

\vspace{0.8cm}

\noindent \textbf{Herramientas utilizadas:} Claude y Codex (principalmente para el desarrollo de código, apoyo en redacción, construcción de diagramas y análisis de resultados).

\vspace{0.6cm}

\noindent \textbf{Reflexión sobre el uso de IA:}

El uso de herramientas de inteligencia artificial durante el proyecto tuvo un car\'acter principalmente instrumental y de apoyo al aprendizaje. La IA se utiliz\'o como una herramienta de consulta para comprender mejor distintas librer\'ias de Python, revisar formas correctas de implementar procedimientos de optimizaci\'on y simulaci\'on, y resolver dudas puntuales sobre estructuras de c\'odigo, manejo de datos y visualizaci\'on de resultados. En este sentido, su aporte fue similar al de una fuente de asistencia t\'ecnica interactiva, que permiti\'o acelerar la b\'usqueda de alternativas y contrastar posibles enfoques antes de incorporarlos al desarrollo del modelo.

En el desarrollo computacional, la IA fue \'util para apoyar la implementaci\'on de los modelos de Markowitz, Black-Litterman y simulaci\'on Monte Carlo, especialmente en tareas relacionadas con el uso de librer\'ias, organizaci\'on del c\'odigo y correcci\'on de errores. Tambi\'en se emple\'o para revisar la consistencia de partes ya desarrolladas, detectar posibles problemas de implementaci\'on y proponer mejoras de eficiencia. Sin embargo, las decisiones metodol\'ogicas, la selecci\'on de supuestos, la interpretaci\'on de resultados y la evaluaci\'on final de las recomendaciones fueron realizadas por el grupo, verificando que las salidas fueran coherentes con los objetivos del proyecto y con la teor\'ia financiera estudiada.

La experiencia mostr\'o que la IA puede aumentar la productividad cuando se usa de manera cr\'itica, pero tambi\'en requiere validaci\'on constante. Sus respuestas no fueron tomadas como definitivas: cada sugerencia relevante fue revisada, adaptada y contrastada con el funcionamiento del c\'odigo, los resultados obtenidos y los criterios del proyecto. Esto fue especialmente importante en un contexto financiero, donde errores de programaci\'on o interpretaci\'on pueden cambiar de forma significativa las conclusiones. Por ello, el uso responsable de IA no reemplaz\'o el razonamiento del equipo, sino que ayud\'o a mejorar la calidad, eficiencia y claridad del trabajo desarrollado.


\vspace{0.25cm}

\begin{center}
\begin{tabular}{p{7cm}l}
Patricio Andres Acevedo Flores & \rule{5cm}{0.4pt} \\[0.55cm]
Benjamín Alonso Cano Araya & \rule{5cm}{0.4pt} \\[0.55cm]
Matías Benjamín Catalán Contreras & \rule{5cm}{0.4pt} \\[0.55cm]
Gaspar Ten Berg Yunusic & \rule{5cm}{0.4pt} \\[0.55cm]
Diego Benjamín Villalobos Michea & \rule{5cm}{0.4pt}
\end{tabular}
\end{center}
